% Part: first-order-logic
% Chapter: natural-deduction
% Section: rules-and-proofs

\documentclass[../../../include/open-logic-section]{subfiles}

\begin{document}

\iftag{FOL}
      {\olfileid{fol}{ntd}{rul}}
      {\olfileid{pl}{ntd}{rul}}

\olsection{Regels en \usetoken{P}{derivation}}

\begin{explain}
Systemen van natuurlijke deductie zijn bedoeld om nauw aan te sluiten
bij de informele redeneringen in wiskundige bewijzen (vandaar dat zij
enigszins ``natuurlijk'' zijn). Bewijzen in natuurlijke deductie
beginnen met aannamen. Vervolgens worden afleidingsregels toegepast.
Aannamen worden door de afleidingsregels \Intro{\lnot}, \Intro{\lif},
\iftag{FOL}{\Elim{\lor} en \Elim{\lexists}}{ en \Elim{\lor}}
``!!{discharged}''. Ter verduidelijking wordt het label van de
!!{discharged} aanname naast de afleidingsstap geplaatst.
\end{explain}

\begin{defn}[Aanname]
Een \emph{aanname} is elke !!{sentence} die bovenaan een tak staat.
\end{defn}

!!^{derivation}s in natuurlijke deductie zijn bepaalde bomen van
!!{sentence}s. De bovenste !!{sentence}s zijn aannamen. Als
!!a{sentence} onder één, twee of drie andere sequenten staat, moet hij
daaruit correct volgen door een afleidingsregel. De !!{sentence}s
boven de afleidingsstap heten de \emph{premissen}; de !!{sentence}
eronder heet de \emph{conclusie} van de afleidingsstap. De regels komen
in paren: voor elke !!{operator} is er een introductieregel en een
eliminatieregel. Zij voeren !!a{operator} in de conclusie in of
verwijderen !!a{operator} uit een premisse van de regel. Sommige regels
staan toe dat een aanname van een bepaald type wordt
\emph{!!{discharged}}. Om aan te geven welke aanname door welke
afleidingsstap wordt !!{discharged}, kennen we zowel de aanname als de
afleidingsstap een label toe. Dit geven we aan door de aanname als
``$\Discharge{!A}{n}$'' te schrijven.

% Only include this sentence is on of \land, lor, lif or lnot is defined.

\iftag{notprvNot,notprvAnd,notprvOr,notprvIf}{}%
	{Het is gebruikelijk regels te beschouwen voor alle !!{operator}s $\land$, $\lor$, $\lif$, $\lnot$ en $\lfalse$, ook als sommige daarvan gedefinieerd zijn.}

\end{document}
